% =====================================================================
%  Working note: deep unrolling of the Random Walker / Power Watershed
%  Lightly edited record of an exchange, structured for reuse.
%
%  Compile:  pdflatex pw_unrolling_note
%            bibtex   pw_unrolling_note
%            pdflatex pw_unrolling_note  (x2)
%
%  Requires: power_watershed_unrolling.bib in the same directory.
% =====================================================================

\documentclass[11pt,a4paper]{article}

\usepackage[T1]{fontenc}
\usepackage[utf8]{inputenc}
\usepackage{amsmath,amssymb,amsthm}
\usepackage[margin=1in]{geometry}
\usepackage[round]{natbib}
\usepackage{hyperref}
\hypersetup{colorlinks=true,linkcolor=black,citecolor=blue,urlcolor=blue}

% --- structure for the dialectical parts -----------------------------
\newcommand{\claim}[1]{\medskip\noindent\textbf{#1}\quad}
\newenvironment{position}[1]
  {\begin{quote}\noindent\textit{#1:}\ }
  {\end{quote}}

% --- notation --------------------------------------------------------
\newcommand{\Lap}{L}
\newcommand{\LU}{\Lap_{U}}
\newcommand{\xU}{x_{U}}
\newcommand{\xM}{x_{M}}
\newcommand{\we}{w_{e}}

\theoremstyle{definition}
\newtheorem{question}{Question}
\newtheorem{remark}{Remark}

\title{Deep Unrolling of the Random Walker and the\\
       Power Watershed Limit: A Feasibility Note}
\author{Working note\thanks{Record of a technical exchange, restructured
        for circulation. Positions attributed as \emph{HT} are Hugues
        Talbot's; the remainder is the interlocutor's analysis. Section
        \ref{sec:open} lists what remains unverified.}}
\date{\today}

\begin{document}
\maketitle

\begin{abstract}
We assess the feasibility of deep-unrolling the Random Walker (RW)
algorithm \citep{grady2006random} for use as a fast, differentiable
edge-refinement module downstream of a CNN segmenter, and of extending
this to the Power Watershed (PW) \citep{couprie2011power}. The RW itself
unrolls cleanly: it reduces to a sparse SPD linear solve admitting both
explicit unrolling and single-solve adjoint differentiation. The PW is
the $p\to\infty$ limit of the same energy, and we argue that the
obstacle to learning through it is not combinatorial non-differentiability
but a \emph{zero-temperature gradient collapse}: convergence of minimizers
does not entail convergence of sensitivities. This reading turns $p$ into
an annealing temperature and suggests training at finite $p$ while
deploying the exact maximum-spanning-forest implementation at
$p=\infty$.
\end{abstract}

% =====================================================================
\section{The question}
% =====================================================================

\begin{position}{HT}
The target is a fast, differentiable watershed usable as an edge-refinement
step inside a deep segmentation pipeline. The proposed route is the Power
Watershed. First question: can the Random Walker be deep-unrolled?
\end{position}

\noindent
Short answer: yes, and the RW is among the most unrolling-friendly
classical algorithms available. The remainder of Section \ref{sec:rw}
makes this precise; Sections \ref{sec:correction}--\ref{sec:temperature}
address the extension to PW, which is where the difficulty actually lies.

% =====================================================================
\section{The Random Walker unrolls cleanly}
\label{sec:rw}
% =====================================================================

Write the graph Laplacian in block form with respect to seeded ($M$) and
unseeded ($U$) nodes,
\begin{equation}
  \Lap = \begin{pmatrix} \Lap_{M} & B \\ B^{\!\top} & \LU \end{pmatrix},
\end{equation}
so that the RW solution satisfies
\begin{equation}
  \LU\,\xU \;=\; -B^{\!\top}\xM .
  \label{eq:rw}
\end{equation}
The matrix $\LU$ is sparse, symmetric and positive definite. Unrolling
here is not unrolling a nonsmooth or combinatorial procedure --- it is
unrolling a linear solve, the most favourable possible setting.

\subsection{Route A: explicit unrolling}

$K$ steps of Jacobi iteration read
\begin{equation}
  x^{(k+1)} \;=\; D^{-1} W x^{(k)}, \qquad
  x^{(k+1)}\big|_{M} \;=\; \xM ,
\end{equation}
i.e.\ a normalised graph convolution with Dirichlet boundary conditions.
The unrolled solver \emph{is} a message-passing network. Conjugate
gradients unrolls equally well and provides a natural slot for a learned
preconditioner. Memory cost is $O(KN)$; see \citet{gregor2010learning}
and \citet{monga2021algorithm} for the general pattern.

\subsection{Route B: implicit differentiation}

Differentiating the optimality condition \eqref{eq:rw} avoids the tape
entirely. For a loss $\ell(\xU)$, define the adjoint
\begin{equation}
  \lambda \;=\; \LU^{-1}\,\frac{\partial \ell}{\partial \xU},
  \qquad \lambda\big|_{M} = 0 ,
\end{equation}
and then for each edge $e=\{u,v\}$,
\begin{equation}
  \boxed{\;
  \frac{\partial \ell}{\partial \we}
  \;=\; -\,(x_u - x_v)\,(\lambda_u - \lambda_v).
  \;}
  \label{eq:adjoint}
\end{equation}
A single adjoint solve --- \emph{with the same matrix}, so the
factorisation or preconditioner is reused --- yields gradients with
respect to every edge weight simultaneously, at memory cost independent
of iteration count \citep{amos2017optnet,bai2019deep}. The derivation is
given in Appendix \ref{app:adjoint}.

\begin{remark}
The product structure of \eqref{eq:adjoint} is itself the argument for
the intended application. In the interior of a segment $x_u \approx x_v$,
so the learning signal concentrates automatically at boundaries. For an
edge-refinement module this is precisely the desired inductive bias.
\end{remark}

\subsection{Prior art}

\citet{cerrone2019learned} implement exactly this: a CNN predicts the
edge weights of the RW graph, interpretable as learning context-dependent
diffusivities for a linear diffusion process. They derive the exact
gradient and, finding the dense Jacobian prohibitive, propose sparse
sampling schemes that remain competitive. PyTorch code is
available \citep{cerrone2019code}. Any new contribution should therefore
be located in the unrolled/truncated solver and the refinement framing,
not in differentiability per se.

\subsection{The truncation caveat}

Truncating at $K$ iterations propagates information $K$ hops. For global
seeded segmentation on a $512^2$ grid this is fatal: one needs
$O(\mathrm{diameter})$ iterations, destroying the speed argument.

For edge refinement it is not. Seeds are the confident interior of the
CNN prediction and the boundary need only move a few pixels, so
$K \sim 5$--$15$ suffices and truncation becomes a design choice rather
than an approximation error. \textbf{This matching of application to the
method's weakness is the strongest single point in the feasibility case.}
If global reach is later required, the standard escape is a learned
multigrid or coarse-to-fine solver rather than more iterations.

\begin{remark}
Prefer the prior-augmented formulation of \citet{grady2005multilabel} ---
adding $\gamma\,\mathrm{diag}(s)$ to $\LU$ with a matching right-hand
side --- over hard seed clamping. It preserves positive definiteness
without hard constraints, makes the seeds differentiable, and provides
the natural port for injecting the CNN's per-pixel logits as unaries.
\end{remark}

% =====================================================================
\section{Correction: PW \emph{is} the RW at $p\to\infty$}
\label{sec:correction}
% =====================================================================

\begin{position}{HT}
The Power Watershed is the Random Walker converging with $p\to\infty$.
The two-phase procedure --- compute a maximum spanning forest, then run
the RW on the plateau of maximal-weight edges --- is a discrete
implementation of that limit. The convergences are the same.
\end{position}

\noindent
This is correct, and an earlier framing of PW as ``a different
algorithm'' was wrong. The unifying energy is
\begin{equation}
  \arg\min_{x}\ \sum_{e=\{u,v\}\in E} \we^{\,p}\,\lvert x_u - x_v\rvert^{q},
  \qquad x\big|_{\text{seeds}} \text{ fixed},
  \label{eq:pw}
\end{equation}
with RW at ($p$ finite, $q=2$) and PW at ($p\to\infty$, $q$ finite)
\citep{couprie2009power,couprie2011power}. At $q=2$ the limit is a
multilabel, scale- and contrast-invariant \emph{unique} global optimum,
obtained in practice in quasi-linear time. The MSF-then-RW-on-plateaus
procedure is an exact finite-time evaluation of that limit, not a
substitute for it. See also
\citet{allene2007links,cousty2009watershed} for the surrounding
equivalences.

% =====================================================================
\section{Where the difficulty actually lives}
\label{sec:sensitivity}
% =====================================================================

The correction stands, and the obstacle survives it in sharper form.

\claim{Convergence of minimizers does not give convergence of sensitivities.}
That $x^{p} \to x^{\infty}$ holds is established. What learning requires
is that $\partial x^{p}/\partial w$ converge to something usable, and
that is a separate limit. As $p$ grows, $x^{\infty}$ becomes locally
constant in $w$ away from changes in plateau structure: perturbing an
edge weight leaves the MSF unchanged, hence leaves the solution
unchanged. Therefore
\begin{equation}
  \frac{\partial x^{\infty}}{\partial w} \;\longrightarrow\; 0
  \quad \text{a.e.,}
\end{equation}
and is undefined on the measure-zero set where the MSF combinatorics
flip. Gradients do not fail to \emph{exist}; they fail to \emph{inform}.

% =====================================================================
\section{$p$ as an inverse temperature}
\label{sec:temperature}
% =====================================================================

The correction improves the plan rather than rescuing it. In the
Probabilistic Watershed formulation \citep{fitasanmartin2019probabilistic},
edge weights are written
\begin{equation}
  w(e) = \exp\!\big(-\mu\, c(e)\big),
\end{equation}
with $\mu$ setting the entropy of the Gibbs distribution over $2$-forests;
raising weights to a power rescales $\mu$. Hence
\begin{equation}
  w(e)^{p} \;=\; \exp\!\big(-p\,\mu\, c(e)\big),
\end{equation}
so $p$ is literally an inverse temperature, and $p\to\infty$ retains only
the minimum-cost spanning forests. The degeneracy of Section
\ref{sec:sensitivity} is then the standard zero-temperature gradient
collapse, structurally identical to
$\mathrm{softmax}\to\arg\max$ \citep{jang2017categorical} and to
entropic optimal transport approaching the exact problem
\citep{cuturi2013sinkhorn}.

\subsection{Proposed scheme}

\begin{enumerate}
  \item \textbf{Train at finite $p$, annealed upward.} Finite-$p$ RW is
    not an approximation of a different method; it is the same
    one-parameter family at nonzero temperature. Unlike learned-watershed
    approaches \citep{wolf2017learned} that train a surrogate and deploy a
    combinatorial procedure, there is no train/test object mismatch: one
    anneals along a path already proven to terminate at the deployed
    algorithm.
  \item \textbf{Deploy at $p=\infty$} with the MSF-plus-plateau
    implementation, retaining the quasi-linear runtime that motivated the
    exercise.
\end{enumerate}

\noindent
\citet{agustsson2017soft} document both annealing failure modes: too fast
and gradients vanish before learning occurs; too slow and the network
inverts the soft assignment by driving weights large. The second mode is
concrete here, since the CNN \emph{predicts} $w$ and can counteract $p$ by
shrinking its output dynamic range. Constraining the output scale is
advisable, and the contrast invariance of PW at $q=2$ supplies a
principled normalisation.

\subsection{The central tension, and a candidate resolution}

Annealing $p$ upward multiplies weight ratios, so
\begin{equation}
  \kappa(\LU) \;\sim\; \exp\!\big(p \cdot \text{(cost spread)}\big),
\end{equation}
and CG iteration counts scale as $\sqrt{\kappa}$. The truncated solver of
Section \ref{sec:rw} therefore degrades exactly as annealing proceeds:
the two ideas fight. \textbf{Measuring the usable $p$ range against the
iteration budget at target resolution is the first experiment.}

A hybrid may dissolve the tension --- namely the existing two-phase
implementation with its second phase warmed:
\begin{center}
  \emph{MSF $+$ finite-$p$ RW restricted to plateaus.}
\end{center}
The forward pass stays near-exact and near-quasi-linear. The gradient is
then supported exactly on plateau edges, i.e.\ on those edges where the
MSF is genuinely ambiguous and the decision is actually being made.
Plateaus are small, so the ill-conditioned solve is small too.

Supporting evidence: \citet{cerrone2019learned} found that arbitrary
sparse sampling of the RW gradient remained competitive. If arbitrary
sparsification suffices, a \emph{principled} sparsification supported on
the ambiguity set should do at least as well.

% =====================================================================
\section{Novelty assessment}
% =====================================================================

\citet{wolf2017learned} attack the combinatorial phase directly, learning
inside Prim's algorithm with structured training on sparse edge sets, and
note that the resulting training set is far smaller than what
unstructured learning affords. \citet{wolf2018mutex} pursue a different
parameter-free partitioning route. Using $p$ itself as the annealing
temperature --- exploiting the PW framework's own parameterisation as the
relaxation path --- appears unoccupied. That the relaxation is
\emph{internal to the framework} rather than bolted on is the
publishable element.

% =====================================================================
\section{Open questions and next steps}
\label{sec:open}
% =====================================================================

\begin{enumerate}
  \item Is the convergence $x^{p}\to x^{\infty}$ uniform enough to bound
    the surrogate gradient bias? A quantitative rate would convert the
    scheme from heuristic to justified.
  \item How does plateau structure behave under CNN-predicted weights?
    Learned weights may produce far fewer exact ties than hand-designed
    ones, shrinking plateaus and starving the gradient. An explicit tie
    tolerance $\varepsilon$ may be required, with its own schedule.
  \item Conditioning study: usable $p$ range versus CG/unrolled iteration
    budget at target resolution.
  \item Plateau-restricted adjoint: derivation and implementation.
  \item Annealing schedule for $p$, jointly with any $\varepsilon$
    schedule and the CNN output normalisation.
\end{enumerate}

% =====================================================================
\appendix
\section{Derivation of the adjoint gradient}
\label{app:adjoint}
% =====================================================================

Take the Dirichlet energy $E(x,w) = \sum_{e=\{u,v\}} \we (x_u-x_v)^2$
with $x|_{M}$ fixed. Stationarity in $\xU$ gives
$\LU \xU + B^{\!\top}\xM = 0$, which is \eqref{eq:rw}. Differentiating
with respect to a single weight $\we$,
\begin{equation}
  \LU \frac{\partial \xU}{\partial \we}
  + \Big[\tfrac{\partial \Lap}{\partial \we} x\Big]_{U} = 0
  \quad\Longrightarrow\quad
  \frac{\partial \xU}{\partial \we}
  = -\,\LU^{-1}\Big[\tfrac{\partial \Lap}{\partial \we} x\Big]_{U}.
\end{equation}
Since $\partial \Lap/\partial \we$ is the elementary edge Laplacian,
the vector $(\partial \Lap/\partial \we)\,x$ has entry $(x_u - x_v)$ at
$u$, entry $(x_v - x_u)$ at $v$, and zero elsewhere. Using symmetry of
$\LU$ and setting $\lambda = \LU^{-1}\partial\ell/\partial\xU$ (extended
by zero on seeds),
\begin{equation}
  \frac{\partial \ell}{\partial \we}
  = \Big(\frac{\partial \ell}{\partial \xU}\Big)^{\!\top}
    \frac{\partial \xU}{\partial \we}
  = -\,\lambda^{\!\top}\Big[\tfrac{\partial \Lap}{\partial \we} x\Big]_{U}
  = -\,(x_u - x_v)(\lambda_u - \lambda_v),
\end{equation}
which is \eqref{eq:adjoint}. For the powered weights of \eqref{eq:pw} the
chain rule contributes an extra factor $p\,\we^{\,p-1}$, whose growth in
$p$ is the algebraic counterpart of the conditioning blow-up.

% =====================================================================
\bibliographystyle{plainnat}
\bibliography{power_watershed_unrolling}
% =====================================================================

\end{document}
